\section*{Chapter \ref{ch:s1:news-and-events} --- News and Events}

\begin{solution}{exo:s1:news-and-events:1}
$e^{-0.1/0.2} = e^{-0.5} = 0.607$ at 100 milliseconds; $e^{-1/0.2} = e^{-5} = 0.0067$ at one second, less than one percent.
\end{solution}

\begin{solution}{exo:s1:news-and-events:2}
$1/(1 + 0.03 \times 39) = 0.46$, so the immediate share is $0.5 + 0.5 \times 0.46 = 0.73$ and 27\% drifts. At the median attention, $0.5 + 0.5 \times 0.42
= 0.71$.
\end{solution}

\begin{solution}{exo:s1:news-and-events:3}
While the stock is halted no one can trade it; the immediate move happens at the reopening auction, where every order meets at one price. There is no
window in which a fast trader can trade at the old price.
\end{solution}

\begin{solution}{exo:s1:news-and-events:4}
The correlation between tone and reading is $1/\sqrt{1 + 1.5^2} = 0.555$, so the signs agree with probability $\tfrac12 + \arcsin(0.555)/\pi = 0.687$;
the simulation measures 68.9\%.
\end{solution}

\begin{solution}{exo:s1:news-and-events:5}
After a second the immediate part is almost entirely gone ($e^{-5}$ of it remains) and the drift has not started: it comes in over the following days.
Between one second and the close nothing moves in the model.
\end{solution}

\begin{solution}{exo:s1:news-and-events:6}
Each split holds about half the items, so its daily P\&L averages over half as many independent positions and its noise is larger. The whole book
diversifies over both halves; its drift per item (25.4 basis points) lies between theirs.
\end{solution}

\begin{solution}{exo:s1:news-and-events:7}
The Sharpe ratio rises to 7.6 and the return to 17.4\% a year, with 45.7 basis points of drift per item before costs instead of 25.4. The difference is the
value of the text model. The trader keeps $25.4/45.7 = 56\%$ of the drift, the correlation between reading and tone (0.555): more than the
$2 \times 0.689 - 1 = 38\%$ that the hit rate alone suggests, because the larger moves are read correctly more often.
\end{solution}

\begin{solution}{exo:s1:news-and-events:8}
A website's publication time is not when the story reached the market: the wire and machine-readable feeds carry it earlier, and machines trade it within
a second. A backtest at the publication minute may trade at prices that already moved, or, if the site's time is earlier than its update, before the
news existed. Use the feed's own timestamps, as received, and the prices after them.
\end{solution}

\begin{solution}{pb:s1:news-and-events:1}
\begin{enumerate}
\item A structured news feed for programs (timestamp, securities, category, scores); the interval from the item plus the latency to the end of the holding
period.
\item Tag the item with its securities; score its direction and strength; decide when to act.
\item High pessimism predicts downward pressure on market prices followed by a reversion to fundamentals, and extreme pessimism high volume.
\item The fraction of negative words forecasts low earnings; prices briefly underreact, most for stories about fundamentals.
\item The reaction can be prepared, which rewards speed; distraction from other announcements is predictable.
\item Weaker immediate reactions and stronger drift when many firms announce the same day; substantial alphas.
\item A stop in trading pending or after news, or after a large move; T1 news pending, T2 news released, LUDP volatility pause.
\item Speed in the first second is worth nothing; the move happens at the reopening auction.
\item 1,000 stocks for eight years; tone with a 2\% standard deviation; attention $1/(1+0.03(k-1))$; an immediate share from a half to one, with a time
constant of 0.2 seconds; the rest over five days; halts above 5\%.
\item 82,351 items, 40.8 a day, 1.18\% halted.
\item 96.7\%, 70.1\%, 29.2\% and 28.7\%.
\item Below a second the machines take the immediate part; above, only the drift remains.
\item \textbf{Named result.} The share of the move still ahead is 96.7\% at 1 millisecond, 93.7\% at 10, 70.1\% at 100, 29.2\% at one second and 28.7\% at
the next close; a daily trader with an imperfect text model captures 25.4 basis points per item, a Sharpe ratio of 3.24 after costs (1.80 on quiet days,
2.74 on busy ones).
\item Attention is spread over more items, so less of each move happens at once: 29.5 against 20.8 basis points of drift per item.
\item 44\% of the drift (it keeps 25.4 of 45.7 basis points), and the Sharpe ratio falls from 7.6 with a perfect reading to 3.24.
\item The model plants the underreaction cleanly and makes the reading's errors independent; real drift is basis points per story and costs absorb much of
it.
\item Feed timestamps as received, prices after them, point-in-time tone scores and costs on every story.
\item Headline sentiment when traded in the first second; the post-halt reopening for its auction data.
\item Chapter~7 traded the drift after earnings in an \omterm{def:s1:earnings:event}{event study}; this chapter measures how much of any item's move is left for each speed.
\item It buys the immediate part of the move; slower traders share what attention leaves behind.
\end{enumerate}
\end{solution}

\begin{solution}{iq:s1:news-and-events:1}
Map each item to its securities and keep the relevant, novel ones; score its tone with a dictionary or a trained model; decide the reaction window and
the horizon; test the signal's IC by horizon on point-in-time data; size it after costs.
\end{solution}

\begin{solution}{iq:s1:news-and-events:2}
The time the feed delivered each item to the firm, against the article's or vendor's stated time; the time of the prices used for entry; and whether
tone scores or tags were revised after the fact.
\end{solution}

\begin{solution}{iq:s1:news-and-events:3}
Read the news when released, estimate the reopening price from the auction's indicative price and imbalance, decide whether the position should be kept,
cut or hedged in related instruments, and place orders in the reopening auction rather than chasing the first trades after it.
\end{solution}

\begin{solution}{iq:s1:news-and-events:4}
A direct feed decoded in a process next to the trading engine, pre-computed mappings from entities to securities, a scoring model small enough to run in
microseconds, pre-staged orders with risk checks done in advance, co-location with the exchange, and timestamps at every step to measure the latency.
\end{solution}

\begin{solution}{iq:s1:news-and-events:5}
Investors have limited attention: with many announcements the same day each one gets less, so prices react less at once and drift more afterwards, as
Hirshleifer, Lim and Teoh found.
\end{solution}

\begin{solution}{iq:s1:news-and-events:6}
At latency $L$ the immediate part still ahead is $a e^{-L/\tau}$ and the drift $1 - a$ is untouched, so the share still ahead is $a e^{-L/\tau} + 1 - a$.
Half of the immediate part is gone at $e^{-L/\tau} = \tfrac12$, $L = \tau \ln 2$: 0.14 seconds for $\tau = 0.2$.
\end{solution}
